\section{Background and Motivation}

The ns-3 TCP model includes multiple congestion-control algorithms, including Linux Reno, CUBIC, DCTCP, and BBR. This capability builds on a substantial redesign of ns-3 TCP architecture and later model contributions for both loss-based and model-based congestion control~\cite{Casoni2016NextGenerationTcp,Nguyen2016TcpVariants,Jain2018Bbr}. The traffic-control module introduces a Linux-like queue-disc layer in ns-3, enabling studies of CoDel, FQ-CoDel, PIE, RED, and related variants~\cite{Imputato2016TrafficControl,Shravya2016Pie}. These models make ns-3 a natural platform for transport and queue-management studies, but they also introduce a large configuration space.

TCP/AQM comparisons are particularly sensitive to hidden assumptions. ECN-enabled queues change whether congestion is signaled by marking or dropping. RTT changes the time required for a flow to reach steady state. Queue-disc defaults affect target delay, marking thresholds, and drop behavior. Some queue disciplines include randomness, so single-run results may overstate a conclusion. These sensitivities are not flaws in ns-3; they are reasons to make experiment design explicit and reproducible.

The starting point for this work is the in-tree TCP validation program. That example provides a dumbbell topology, single-flow and optional second-flow TCP traffic, ping RTT measurement, bottleneck queue tracing, and validation cases for selected CUBIC and DCTCP behavior. Our benchmark preserves that topology as an anchor but moves paper-specific generalization into a new example so the original validation example remains unchanged.
